\usepackage{ragged2e}
\usepackage{etoolbox}

\apptocmd{\frame}{}{\justifying}{}

\setbeamertemplate{enumerate items}[default]
\setbeamertemplate{itemize items}[circle]
\setbeamersize{text margin left=2.25em, text margin right=2.25em}

\usepackage{orcidlink}
\usepackage{amsmath}
\DeclareMathOperator*{\argmax}{argmax}

% A slide is only ~46 characters wide at the default font size, and a
% code line that overruns it does not wrap -- it runs off the right edge
% of the frame and off the page. Quarto builds highlighted code blocks
% as fancyvrb Verbatim environments, so \fvset makes them wrap instead.
% breakanywhere is needed because the long lines are usually paths or
% LaTeX commands with no spaces to break at. This is a safety net, not
% a licence: wrapped code still reads badly on a slide, so keep code
% samples short.
%
% fontsize sets the code font for highlighted blocks. \footnotesize fits
% roughly 65 characters across a 4:3 frame (\small fits ~55,
% \scriptsize ~80), which covers most realistic code lines without
% wrapping and still reads on a projector. Inline `code` is \texttt, not
% fancyvrb, so it keeps the body size -- that is deliberate, a smaller
% inline size looks broken mid-sentence.
\usepackage{fvextra}
\fvset{fontsize=\footnotesize, breaklines=true, breakanywhere=true}

% ---------------------------------------------------------------------
% List spacing.
%
% \qkitlistsep is the space BETWEEN items. At the TOP level \topsep is
% zero, so the space around a list equals the space between two
% paragraphs. That keeps one rhythm on a slide: blocks (paragraph, list)
% are separated by the paragraph skip, items inside a list by a smaller
% \qkitlistsep, and display equations by \qkitdisplayskip (below)
% regardless of what sits next to them. Top-level \topsep is the length
% that property hangs on -- give it a value and a display equation
% adjacent to a list picks up that value on top of \qkitdisplayskip,
% while one adjacent to a paragraph does not.
%
% NESTED lists do take \topsep, because they are not the boundary
% between two blocks: a nested itemize opens flush under the text of its
% parent item, and at zero \topsep it sat one line below it, tighter
% than the gap between its own items. \qkitlistsep there makes entering
% a nested list cost exactly what stepping to the next item costs.
% Only the opening gap moves -- returning from the nested list to the
% next parent item is set by the parent's \itemsep and is unaffected.
%
% \tightlist is why this has to be done here. Pandoc emits it right
% after \begin{itemize} / \begin{enumerate} for every list whose source
% has no blank lines between items -- which is nearly every list -- and
% its default definition sets \itemsep to 0pt. Anything that adjusts
% \itemsep on the way INTO the list (the \addtolength{\itemsep} hack
% this replaced) is therefore undone a token later, which is why that
% hack had no visible effect. Quarto declares \tightlist with
% \providecommand ahead of this preamble, so redefining it here wins;
% the \providecommand below is a no-op guard in case a future Quarto
% stops emitting it at all.
%
% \@listi / \@listii / \@listiii carry the same \itemsep so that a
% NON-tight list -- one with blank lines between its items, which never
% sees \tightlist -- is spaced identically. They otherwise keep beamer's
% own values from beamerbaselocalstructure.sty.
% ---------------------------------------------------------------------
\newlength{\qkitlistsep}
\setlength{\qkitlistsep}{4pt}
\providecommand{\tightlist}{}
\renewcommand{\tightlist}{%
  \setlength{\itemsep}{\qkitlistsep}%
  \setlength{\parskip}{0pt}%
}
\makeatletter
\def\@listi{\leftmargin\leftmargini
            \topsep    \z@ \@plus\p@
            \parsep    0\p@
            \itemsep   \qkitlistsep}
\let\@listI\@listi
\def\@listii{\leftmargin\leftmarginii
             \topsep    \qkitlistsep
             \parsep    0\p@   \@plus\p@
             \itemsep   \qkitlistsep}
\def\@listiii{\leftmargin\leftmarginiii
              \topsep    \qkitlistsep
              \parsep    0\p@   \@plus\p@
              \itemsep   \qkitlistsep}
\makeatother

% ---------------------------------------------------------------------
% Display-equation spacing: the same gap above and below, and the same
% gap whether the display sits next to a paragraph, an itemize, or an
% enumerate.
%
% Two things are needed. First the four display skips are pinned to one
% value -- the short-skip variants apply when the preceding line is
% short, and by default \abovedisplayshortskip is 0pt, so a display
% after a short line hugs it. Second, and much the larger effect,
% \qkitdisplayfix cancels the empty paragraph above the display: Pandoc
% always emits a display as its own block, so LaTeX meets
% \begin{equation} in VERTICAL mode, opens a paragraph, finds it empty,
% and contributes a full \baselineskip that nothing below the display
% matches. The guards mean a display that genuinely continues a
% paragraph (hand-written LaTeX, no blank line) is left alone, and
% nothing is subtracted at the top of a frame where there is no
% interline glue to cancel.
%
% Measured on a 4:3 frame: 30.3pt above and 30.2pt below, identical
% after a paragraph, an itemize and an enumerate.
% ---------------------------------------------------------------------
\newlength{\qkitdisplayskip}
\setlength{\qkitdisplayskip}{10pt}
\newcommand{\qkitdisplayskips}{%
  \setlength{\abovedisplayskip}{\qkitdisplayskip}%
  \setlength{\belowdisplayskip}{\qkitdisplayskip}%
  \setlength{\abovedisplayshortskip}{\qkitdisplayskip}%
  \setlength{\belowdisplayshortskip}{\qkitdisplayskip}%
}
\makeatletter
\newcommand{\qkitdisplayfix}{%
  \qkitdisplayskips
  \ifvmode\ifdim\prevdepth>-\@m pt \vskip-\baselineskip\fi\fi
}
\makeatother
\qkitdisplayskips
\AtBeginDocument{\qkitdisplayskips}
\everydisplay=\expandafter{\the\everydisplay \qkitdisplayskips}
% amsmath defines \[ ... \] as the equation* environment, so the
% equation* hook also covers unlabelled `$$ ... $$` displays.
\BeforeBeginEnvironment{equation}{\qkitdisplayfix}
\BeforeBeginEnvironment{equation*}{\qkitdisplayfix}
\BeforeBeginEnvironment{align}{\qkitdisplayfix}
\BeforeBeginEnvironment{align*}{\qkitdisplayfix}
\BeforeBeginEnvironment{gather}{\qkitdisplayfix}
\BeforeBeginEnvironment{gather*}{\qkitdisplayfix}
\BeforeBeginEnvironment{multline}{\qkitdisplayfix}
\BeforeBeginEnvironment{multline*}{\qkitdisplayfix}

\let\olditem\item
\renewcommand\item{\olditem\justifying}

\definecolor{blue}{RGB}{0,114,178}
\definecolor{red}{RGB}{213,94,0}
\definecolor{yellow}{RGB}{240,228,66}
\definecolor{green}{RGB}{0,158,115}

\definecolor{MyBackground}{RGB}{255,253,218}

\setbeamercolor{section in toc}{fg=blue}
\setbeamercolor{subsection in toc}{fg=red}

\setbeamercolor{frametitle}{fg=blue}
\setbeamercolor{title}{fg=blue}

\setbeamertemplate{footline}[frame number]
\setbeamertemplate{navigation symbols}{}
\setbeamertemplate{itemize items}{-}

\setbeamercolor{itemize item}{fg=blue}
\setbeamercolor{itemize subitem}{fg=blue}
\setbeamercolor{enumerate item}{fg=blue}
\setbeamercolor{enumerate subitem}{fg=blue}
\setbeamercolor{button}{bg=MyBackground,fg=blue}

\newcounter{daggerfootnote}
\newcommand*{\daggerfootnote}[1]{%
    \setcounter{daggerfootnote}{\value{footnote}}%
    \renewcommand*{\thefootnote}{\fnsymbol{footnote}}%
    \footnote[2]{#1}%
    \setcounter{footnote}{\value{daggerfootnote}}%
    \renewcommand*{\thefootnote}{\arabic{footnote}}%
}

\setbeamertemplate{date}{
  \vspace{0.8em}
  \insertdate
}

\AtBeginSection{
    \begingroup
    \setbeamercolor{background canvas}{bg=yellow}
    \begin{frame}
        \begin{centering}
            \usebeamerfont{section title}\huge\color{black}\insertsection\par
        \end{centering}
    \end{frame}
    \endgroup
}

% allow use of internal LaTeX commands (@)
\makeatletter

\newcommand{\BLButtons}{} % store buttons for the current slide

% add a button to the bottom-left (can be called multiple times)
\newcommand{\BLButton}[1]{%
  \g@addto@macro\BLButtons{#1\hspace{0.6em}}%
}

% clear buttons at the start of each frame
\AtBeginEnvironment{frame}{\gdef\BLButtons{}}

% custom footline
\setbeamertemplate{footline}{
  \begin{beamercolorbox}[
    wd=\paperwidth,
    ht=2.8ex,
    dp=1.2ex,
    leftskip=1em,
    rightskip=1em
  ]{footline}

    \BLButtons        % buttons on the left
    \hfill            % push content apart
    \insertframenumber{} / \inserttotalframenumber % frame numbers

  \end{beamercolorbox}
}

% end internal command access
\makeatother
